% ============================================================================
% ALPHA ASYMMETRY IN FOREIGN EXCHANGE MARKETS
% Farzulla Research Working Paper
% ============================================================================

\documentclass[letterpaper,11pt]{article}

% ============================================================================
% PACKAGE IMPORTS
% ============================================================================

\usepackage[utf8]{inputenc}
\usepackage[T1]{fontenc}
\usepackage{mathpazo}  % Palatino font (matches ROM/arXiv style)
\usepackage[letterpaper, margin=1in]{geometry}
\usepackage{setspace}
\setstretch{1.1}
\usepackage{graphicx}
\usepackage{float}
\usepackage{booktabs}
\usepackage{array}
\usepackage{longtable}
\usepackage{placeins}
\usepackage{amsmath}
\usepackage{amssymb}
\usepackage{amsthm}
\usepackage[round,authoryear]{natbib}
\bibliographystyle{plainnat}

% Colors
\usepackage{xcolor}
\definecolor{farzullaburgundy}{RGB}{128,0,32}
\definecolor{zenodoblue}{RGB}{0,123,255}
\definecolor{orcidgreen}{RGB}{166,206,57}

\usepackage{tcolorbox}
\usepackage{tikz}
\usepackage{scalerel}

% Hyperlinks
\usepackage{url}
\usepackage[colorlinks=true,
            linkcolor=farzullaburgundy,
            citecolor=farzullaburgundy,
            urlcolor=farzullaburgundy,
            breaklinks=true,
            pdftitle={Alpha Asymmetry in Foreign Exchange Markets},
            pdfauthor={Murad Farzulla},
            pdfkeywords={alpha asymmetry, forex, skewness, kurtosis, trading strategies}]{hyperref}

\def\UrlBreaks{\do\/\do-\do_}
\expandafter\def\expandafter\UrlBreaks\expandafter{\UrlBreaks\do\a\do\b\do\c\do\d\do\e\do\f\do\g\do\h\do\i\do\j\do\k\do\l\do\m\do\n\do\o\do\p\do\q\do\r\do\s\do\t\do\u\do\v\do\w\do\x\do\y\do\z}

% Section formatting
\usepackage{titlesec}
\titlespacing*{\section}{0pt}{1.5ex plus 0.5ex minus 0.2ex}{1ex plus 0.2ex}
\titlespacing*{\subsection}{0pt}{1.2ex plus 0.4ex minus 0.2ex}{0.8ex plus 0.2ex}
\titleformat{\section}{\normalfont\large\bfseries\color{farzullaburgundy}}{\thesection}{0.5em}{}
\titleformat{\subsection}{\normalfont\normalsize\bfseries\color{farzullaburgundy}}{\thesubsection}{0.5em}{}
\titleformat{\subsubsection}{\normalfont\small\bfseries\color{farzullaburgundy}}{\thesubsubsection}{0.5em}{}

\usepackage{fancyhdr}

% ============================================================================
% CUSTOM COMMANDS
% ============================================================================

% ORCID icon
\newcommand{\orcidicon}{\scalerel*{
    \begin{tikzpicture}[x=3ex,y=3ex]
    \draw[fill=orcidgreen] (0,0) circle (0.5);
    \draw[white,line width=0.08ex] (0,0.15) -- (0,0.3);
    \draw[white,line width=0.08ex] (-0.15,0) -- (-0.3,0);
    \end{tikzpicture}
}{\textrm{I}}}

% Farzulla Research logo (TikZ version)
\newcommand{\farzullalogo}{%
    \href{https://farzulla.org}{%
        \scalerel*{%
            \begin{tikzpicture}[x=1ex,y=1ex]
            \fill[farzullaburgundy] (0,0) rectangle (1,1.5);
            \fill[white] (0.2,0.2) rectangle (0.8,0.5);
            \fill[white] (0.2,0.7) rectangle (0.6,0.9);
            \fill[white] (0.2,1.1) rectangle (0.8,1.3);
            \end{tikzpicture}%
        }{\textrm{F}}%
    }%
}

% Zenodo logo (TikZ version)
\newcommand{\zenodologo}{%
    \href{https://zenodo.org}{%
        \scalerel*{%
            \begin{tikzpicture}[x=1ex,y=1ex]
            \fill[zenodoblue] (0,0) circle (0.75);
            \fill[white] (0,0) circle (0.5);
            \fill[zenodoblue] (0,0) circle (0.25);
            \end{tikzpicture}%
        }{\textrm{Z}}%
    }%
}

% Version info
\newcommand{\paperver}{2.0.0}
\newcommand{\paperdate}{January 2026}
\newcommand{\paperdoi}{10.5281/zenodo.17918374}
\newcommand{\shorttitle}{Alpha Asymmetry in FX Markets}

% Metadata box
\newcommand{\metadatabox}[1]{%
\begin{tcolorbox}[
    colback=gray!5,
    colframe=farzullaburgundy,
    title=Publication Metadata,
    fonttitle=\bfseries,
    coltitle=white,
    colbacktitle=farzullaburgundy,
    width=\columnwidth,
    arc=2mm,
    boxrule=0.8pt
]
\small
\textbf{DOI:} \href{https://doi.org/#1}{\texttt{#1}}\\
\textbf{Version:} \paperver\\
\textbf{Date:} \paperdate\\
\textbf{License:} CC-BY-4.0
\end{tcolorbox}
}

% ============================================================================
% HEADER/FOOTER
% ============================================================================

\pagestyle{fancy}
\fancyhf{}
\fancyhead[L]{\small\itshape\shorttitle}
\fancyhead[R]{\small\thepage}
\fancyfoot[C]{}
\renewcommand{\headrulewidth}{0.4pt}
\renewcommand{\footrulewidth}{0pt}

\fancypagestyle{firstpage}{
  \fancyhf{}
  \fancyfoot[C]{\small\thepage}
  \renewcommand{\headrulewidth}{0pt}
  \renewcommand{\footrulewidth}{0pt}
}

% ============================================================================
% DOCUMENT
% ============================================================================

\begin{document}

% Single-column throughout (standardized formatting)
\thispagestyle{firstpage}

% Title block
\begin{center}
{\Large\bfseries Alpha Asymmetry in Foreign Exchange Markets}\\[0.5em]
{\large\itshape An Investigation of Exploitability}\\[1em]
{\bfseries Murad Farzulla}\textsuperscript{1,2}\\[0.5em]
{\small\itshape \textsuperscript{1}King's College London, \textsuperscript{2}\href{https://dissensus.ai}{Dissensus AI}}\\[0.3em]
{\small \paperdate}\\[0.5em]
{\small Correspondence: \href{mailto:murad@dissensus.ai}{murad@dissensus.ai}}
\end{center}

% Abstract
\begin{abstract}
\noindent This paper investigates whether distributional asymmetries in foreign
exchange alpha signals represent exploitable market inefficiencies. Using EUR/JPY
data spanning November 2015--August 2025 (510 weekly observations), we document
statistically significant departures from normality across five alpha types, with
pronounced right-skewness in tail alpha (5.05) and momentum signals (2.12).
However, we find that these asymmetries do not translate to economically
significant trading profits. The GPD shape parameter is not significantly
different from zero ($\xi = -0.23$, 95\% CI: $[-1.79, 0.24]$), indicating
asymmetry arises from outlier frequency rather than heavy tails. Strategy
returns include zero in confidence intervals after HAC correction; cross-market
validation fails for equities and commodities; and transaction costs eliminate
the modest gross edge. We conclude that alpha signal asymmetry, while
statistically detectable, does not constitute an exploitable market inefficiency
in FX markets. These null findings caution against over-interpreting
higher-moment statistics as trading signals without rigorous economic validation.

\vspace{0.5em}
\noindent\textbf{Keywords:} null result, alpha asymmetry, foreign exchange, skewness, market efficiency, extreme value theory

\vspace{0.3em}
\noindent\textbf{JEL Codes:} G11, G14, G15, C58
\end{abstract}

% ============================================================================
% MAIN BODY
% ============================================================================

\vspace{1em}

\section{Introduction}

Financial markets exhibit persistent deviations from the efficient market hypothesis \citep{fama1970efficient}, with alpha signals---measures of risk-adjusted excess returns---displaying systematic patterns that sophisticated traders exploit \citep{jegadeesh1993returns, moskowitz2012time}. While considerable literature examines alpha generation and decay, less attention has been paid to the \textit{distributional properties} of alpha signals themselves. This is surprising given well-documented evidence that asset returns deviate substantially from normality, exhibiting fat tails \citep{mandelbrot1963variation, cont2001empirical} and asymmetric distributions \citep{harvey2000conditional}.

The preference for skewed returns has deep roots in asset pricing theory. \citet{kraus1976skewness} established that investors prefer positive skewness and dislike negative skewness, implying that assets with lottery-like payoffs command lower expected returns. Subsequent work has confirmed that skewness affects both individual asset pricing \citep{harvey2000conditional, bali2011maxing} and portfolio construction \citep{mitton2007equilibrium, brunnermeier2007optimal}.

This paper investigates whether alpha signals exhibit systematic asymmetries in their probability distributions, and whether such asymmetries can inform profitable trading strategies. We test this hypothesis rigorously and report negative findings: while asymmetry is statistically detectable, it does not survive the transition from statistical significance to economic significance. We investigate the following questions:

\begin{enumerate}
    \item Do alpha signals in forex markets deviate significantly from normal distributions? \textit{(Yes---confirmed)}
    \item Do these deviations manifest as exploitable skewness and heavy tails? \textit{(Partially---skewness yes, heavy tails no)}
    \item Do distributional asymmetries persist across market regimes? \textit{(Yes, but inconsistently)}
    \item Do asymmetry-aware strategies outperform benchmarks after costs? \textit{(No---the central null finding)}
\end{enumerate}

Our contribution is methodological honesty: we document a plausible-sounding trading idea that does not survive rigorous testing. Such null results are underreported in quantitative finance \citep{harvey2017presidential}, yet they prevent wasted research effort and capital allocation to spurious patterns. We show precisely where the asymmetry-exploitation thesis fails: not in signal detection (which works), but in the translation from statistical pattern to economic profit.

Our empirical contribution is threefold. First, we provide a systematic characterization of distributional properties across five distinct alpha types in EUR/JPY forex data, contributing to the growing literature on FX market microstructure \citep{king2013market, evans2002order}. Second, we demonstrate that asymmetry-based trading strategies, while generating positive gross returns, do not survive transaction costs and statistical uncertainty---extending work on technical trading rules \citep{brock1992simple, lo2000foundations, neely2014forecasting} by showing where such rules fail. Third, we test cross-market generalizability and find it lacking: asymmetry patterns in SPY and GLD do not translate to profitable strategies, suggesting EUR/JPY results may reflect idiosyncratic microstructure rather than a general principle \citep{asness2013value}.

\section{Methodology}

\subsection{Data and Sample Construction}

We analyze EUR/JPY forex data comprising 510 weekly observations spanning November 2015 to August 2025, sourced from Yahoo Finance (EURJPY=X). The sample period encompasses multiple market regimes including the post-Brexit volatility spike (2016), COVID-19 pandemic shock (2020), and subsequent monetary policy divergence between the Federal Reserve, ECB, and Bank of Japan.

\textbf{Data Frequency and Aggregation.} Alpha signals are computed daily using the specified rolling windows (5-day, 20-day, 60-day), then aggregated to weekly frequency using Friday closing values to align with institutional trading cycles. The 510 weekly observations represent end-of-week snapshots of daily-computed signals, ensuring consistency between high-frequency alpha construction and lower-frequency backtesting. All returns and volatility measures are computed at daily frequency before weekly aggregation.

The dataset includes five pre-computed alpha signals, each capturing distinct aspects of market dynamics:

\begin{itemize}
    \item \textbf{Tail Alpha ($\alpha_{\text{tail}}$)}: Captures extreme price movements using a modified $z$-score of returns beyond the 95th percentile. Formally, $\alpha_{\text{tail},t} = \mathbf{1}_{|r_t| > q_{0.95}} \cdot \text{sgn}(r_t) \cdot |r_t|$, where $q_{0.95}$ denotes the rolling 52-week 95th percentile of absolute returns.

    \item \textbf{Fast Alpha ($\alpha_{\text{fast}}$)}: Short-term momentum signal computed as the 5-day return normalized by 20-day realized volatility: $\alpha_{\text{fast},t} = (P_t - P_{t-5}) / (\sigma_{20,t} \sqrt{5})$.

    \item \textbf{Pricing Alpha ($\alpha_{\text{price}}$)}: Mean-reversion signal measuring deviation from fair value, computed as $\alpha_{\text{price},t} = (P_t - \text{MA}_{60,t}) / \sigma_{60,t}$, where $\text{MA}_{60}$ and $\sigma_{60}$ denote 60-day moving average and standard deviation.

    \item \textbf{Coverage Alpha ($\alpha_{\text{cov}}$)}: Volatility compression ratio measuring regime shifts in realized volatility:
    \begin{equation}
        \alpha_{\text{cov},t} = \frac{\sigma_{20,t}}{\sigma_{20,t-5}} - 1
    \end{equation}
    where $\sigma_{20,t}$ denotes 20-day realized volatility computed from daily close-to-close returns. Data source: Yahoo Finance (EURJPY=X). Values above zero indicate volatility expansion; values below zero indicate compression.

    \item \textbf{Hedge Alpha ($\alpha_{\text{hedge}}$)}: Captures the interaction between dollar correlation and interest rate carry. When EUR/JPY moves in tandem with dollar strength (high positive correlation with DXY), carry trades face amplified risk during dollar rallies. The hedge alpha is positive when: (i) correlation is positive and Japan offers higher rates (favorable for short EUR/JPY hedges); or (ii) correlation is negative and US offers higher rates (favorable for long EUR/JPY hedges). Formally:
    \begin{equation}
        \alpha_{\text{hedge},t} = \rho_{t}^{(\text{EUR/JPY}, \text{DXY})} \times \Delta r_{t}^{(\text{JPY}-\text{USD})}
    \end{equation}
    where $\rho$ is the 20-week rolling correlation between EUR/JPY and DXY (US Dollar Index, DX-Y.NYB from Yahoo Finance), and $\Delta r = r_{\text{JPY}} - r_{\text{USD}}$ is the annualized interest rate differential sourced from FRED (INTDSRJPM193N for Japan, INTDSRUSM193N for US). This alpha captures regimes where directional (correlation) and carry (rate differential) signals align, suggesting exploitable hedging opportunities.
\end{itemize}

For cross-market validation, we construct comparable alpha measures for GBP/USD, SPY (S\&P 500 ETF), and GLD (Gold ETF) using Yahoo Finance data, applying identical transformation methodologies to ensure comparability.

\subsection{Asymmetry Metrics}

We employ four complementary measures to characterize distributional asymmetry, each capturing distinct aspects of non-normality relevant to trading strategy design.

\textbf{Sample Skewness} measures the degree of asymmetry around the mean. For a sample $\{x_1, \ldots, x_n\}$ with sample mean $\bar{x}$ and standard deviation $s$, we compute the adjusted Fisher-Pearson coefficient:
\begin{equation}
    \hat{\gamma}_1 = \frac{n}{(n-1)(n-2)} \sum_{i=1}^{n} \left( \frac{x_i - \bar{x}}{s} \right)^3
\end{equation}
This bias-corrected estimator is consistent under standard regularity conditions. Positive skewness ($\hat{\gamma}_1 > 0$) indicates a right-tailed distribution with more extreme positive values, while negative skewness indicates left-tail heaviness. The standard error of skewness under normality is approximately $\sqrt{6/n} \approx 0.108$ for our sample size.

\textbf{Excess Kurtosis} captures tail heaviness relative to a Gaussian benchmark:
\begin{equation}
    \hat{\gamma}_2 = \frac{n(n+1)}{(n-1)(n-2)(n-3)} \sum_{i=1}^{n} \left( \frac{x_i - \bar{x}}{s} \right)^4 - \frac{3(n-1)^2}{(n-2)(n-3)}
\end{equation}
This estimator subtracts 3 (the kurtosis of a normal distribution) and applies finite-sample bias correction. Values exceeding zero indicate leptokurtic (fat-tailed) distributions; the standard error under normality is approximately $\sqrt{24/n} \approx 0.217$.

\textbf{Asymmetry Index (AI)} quantifies the ratio of upside to downside semi-variance, providing a risk-management-oriented asymmetry measure:
\begin{equation}
    AI = \frac{\text{Var}^+(X)}{\text{Var}^-(X)} = \frac{\sum_{i: x_i > \bar{x}} (x_i - \bar{x})^2 / n^+}{\sum_{i: x_i < \bar{x}} (x_i - \bar{x})^2 / n^-}
\end{equation}
where $n^+$ and $n^-$ denote observations above and below the mean. Values above 1 indicate greater dispersion in positive deviations---relevant for assessing ``lottery ticket'' payoff structures \citep{bali2011maxing}.

\textbf{Positive Observation Ratio (PNR)} measures the unconditional probability of positive realizations:
\begin{equation}
    PNR = \frac{1}{n} \sum_{i=1}^{n} \mathbf{1}_{x_i > 0}
\end{equation}
Under normality with zero mean, $PNR = 0.5$; deviations indicate location-scale asymmetry or non-zero drift.

\subsection{Statistical Tests}

We employ a battery of complementary tests to assess departures from normality, each with distinct power properties against different alternatives.

\textbf{Shapiro-Wilk Test.} The Shapiro-Wilk statistic \citep{shapiro1965analysis} tests the null hypothesis $H_0: X \sim \mathcal{N}(\mu, \sigma^2)$ against general alternatives. The test statistic is:
\begin{equation}
    W = \frac{\left( \sum_{i=1}^{n} a_i x_{(i)} \right)^2}{\sum_{i=1}^{n} (x_i - \bar{x})^2}
\end{equation}
where $x_{(i)}$ are order statistics and $a_i$ are tabulated coefficients. This test has high power against asymmetric and heavy-tailed alternatives for moderate sample sizes.

\textbf{D'Agostino-Pearson Omnibus Test.} Following \citet{dagostino1990suggestion}, we test skewness and kurtosis jointly using the $K^2$ statistic:
\begin{equation}
    K^2 = Z_1(\hat{\gamma}_1)^2 + Z_2(\hat{\gamma}_2)^2 \sim \chi^2_2
\end{equation}
where $Z_1$ and $Z_2$ are normalizing transformations of sample skewness and kurtosis. This test is particularly powerful against asymmetric alternatives.

\textbf{Jarque-Bera Test.} As a robustness check, we employ the Jarque-Bera statistic \citep{jarque1980efficient}:
\begin{equation}
    JB = \frac{n}{6} \left( \hat{\gamma}_1^2 + \frac{(\hat{\gamma}_2)^2}{4} \right) \stackrel{a}{\sim} \chi^2_2
\end{equation}
This Lagrange multiplier test is asymptotically equivalent to the likelihood ratio test for normality.

\textbf{Skewness Significance.} We test $H_0: \gamma_1 = 0$ using the $t$-ratio $t = \hat{\gamma}_1 / \text{SE}(\hat{\gamma}_1)$, where $\text{SE}(\hat{\gamma}_1) \approx \sqrt{6(n-2)/[(n+1)(n+3)]}$. For $n = 510$, this yields a standard error of approximately 0.108.

\textbf{Multiple Testing Correction.} Given five alpha types and multiple test statistics, we apply the Bonferroni correction to control the family-wise error rate at $\alpha = 0.05$, yielding adjusted significance thresholds of $\alpha^* = 0.05/5 = 0.01$ per alpha type.

\subsection{Backtesting Framework}

We implement three trading strategies for comparison, with explicit signal generation and execution rules to ensure reproducibility.

\textbf{Asymmetry Strategy}:

\textit{Signal Generation:}
\begin{itemize}
    \item \textbf{Entry Long}: Rolling skewness of fast alpha (20-week window) exceeds 0.75 AND current fast alpha $> 0$
    \item \textbf{Entry Short}: Rolling skewness of pricing alpha (20-week window) exceeds 0.75 AND current pricing alpha $> 0.5\sigma_{\text{pricing}}$
    \item \textbf{Exit}: Signal reversal (opposite entry condition met) OR 4-week maximum holding period reached
\end{itemize}

\textit{Position Sizing:}
\begin{equation}
    \text{Position} = \max(0.5, \min(2.0, 1 + |AI_t - 1.0|))
\end{equation}
where $AI_t$ is the contemporaneous asymmetry index. This scales positions up when asymmetry is pronounced and down when distributions approach symmetry.

\textit{Execution:}
\begin{itemize}
    \item Entry: Monday open following Friday signal generation
    \item Exit: Friday close or upon signal reversal
    \item Rebalancing: Weekly (end of Friday close)
    \item No leverage; positions bounded to $[0.5, 2.0]$ units
\end{itemize}

\textbf{Momentum Strategy}: Classic trend-following using 20-day moving average crossovers. Long when price crosses above MA(20); short when price crosses below. Exit on reversal.

\textbf{Mean Reversion Strategy}: Contrarian positions when prices deviate more than 2 standard deviations from 20-day mean. Long when $P_t < \text{MA}_{20} - 2\sigma_{20}$; short when $P_t > \text{MA}_{20} + 2\sigma_{20}$. Exit when price returns within 0.5 standard deviations of mean.

Performance metrics include total return, Sharpe ratio, Sortino ratio, maximum drawdown, win rate, and trade count.

\section{Results}

\subsection{Asymmetry Detection}

Table \ref{tab:asymmetry} presents distributional statistics for each alpha type.

\begin{table}[H]
\centering
\caption{Alpha Asymmetry Metrics}
\label{tab:asymmetry}
\small
\begin{tabular}{@{}lcccc@{}}
\toprule
Alpha Type & Skew & Kurt & AI & PNR \\
\midrule
Tail & 5.05 & 47.41 & 1.38 & 3.25\% \\
Fast & 2.12 & 12.70 & 1.35 & 57.52\% \\
Pricing & 1.53 & 5.32 & 1.39 & 45.20\% \\
Coverage & -0.04 & 0.94 & 0.99 & 52.79\% \\
Hedge & -1.45 & 4.43 & 0.71 & 58.76\% \\
\bottomrule
\end{tabular}
\vspace{0.3em}
\footnotesize{Note: Skew = skewness, Kurt = excess kurtosis, AI = asymmetry index, PNR = positive/negative ratio.}
\end{table}

These statistics confirm that alpha signals exhibit pronounced non-normality. However, statistical detectability does not imply economic exploitability---a distinction we examine in the following sections.

Several patterns emerge. First, tail alpha exhibits extreme right-skewness (5.05) and massive kurtosis (47.41), indicating rare but substantial positive outliers---consistent with ``black swan'' events generating outsized returns \citep{taleb2007black}. These distributional properties align with the stylized facts documented in \citet{cont2001empirical}. Second, fast and pricing alphas display moderate right-skewness with heavy tails, suggesting momentum \citep{jegadeesh1993returns} and mean-reversion \citep{debondt1985does} signals cluster asymmetrically. Third, coverage alpha approximates symmetry (skew $\approx$ 0, AI $\approx$ 1), indicating liquidity signals distribute more normally, consistent with microstructure theory \citep{kyle1985continuous, glosten1985bid}. Fourth, hedge alpha exhibits left-skewness (-1.45), suggesting correlation signals produce more extreme negative values---a pattern reminiscent of carry trade crash risk \citep{burnside2011carry}.

\begin{figure*}[t]
\centering
\includegraphics[width=0.95\textwidth]{alpha_asymmetry_analysis.png}
\caption{Alpha Asymmetry Analysis. Distributional properties of five alpha types showing skewness, kurtosis, and asymmetry patterns across EUR/JPY forex data (2015--2025).}
\label{fig:asymmetry}
\end{figure*}

\subsection{Statistical Significance}

Table \ref{tab:tests} presents comprehensive hypothesis test results for each alpha type. We report test statistics and $p$-values for multiple normality and asymmetry tests to ensure robustness.

\begin{table}[H]
\centering
\caption{Comprehensive Statistical Test Results}
\label{tab:tests}
\small
\begin{tabular}{@{}lcccccc@{}}
\toprule
Alpha & $\hat{\gamma}_1$ & $t$-stat & SW & JB & $K^2$ & Normal? \\
\midrule
Tail & 5.05 & 52.6*** & 0.412*** & 8941*** & 2847*** & Rejected \\
Fast & 2.12 & 22.1*** & 0.891*** & 1124*** & 512*** & Rejected \\
Pricing & 1.53 & 15.9*** & 0.924*** & 463*** & 267*** & Rejected \\
Coverage & -0.04 & -0.42 & 0.978*** & 28.4*** & 22.1*** & Rejected \\
Hedge & -1.45 & -15.1*** & 0.932*** & 389*** & 241*** & Rejected \\
\bottomrule
\end{tabular}
\vspace{0.3em}
\footnotesize{Note: $t$-stat tests $H_0: \gamma_1 = 0$; SW = Shapiro-Wilk statistic; JB = Jarque-Bera statistic; $K^2$ = D'Agostino-Pearson omnibus. *** $p < 0.001$ after Bonferroni correction. For robustness, Ljung-Box tests on standardized residuals yield $Q(4) = 8.12$ ($p = 0.087$), indicating modest but non-negligible serial correlation; normality test statistics remain significant after \citet{lobato2001testing} HAC correction for dependent data.}
\end{table}

\textbf{Normality Rejection.} All five alpha types reject the null hypothesis of normality across all three tests at the $p < 0.001$ level, even after Bonferroni correction. The Shapiro-Wilk statistics range from 0.412 (tail alpha, indicating severe departure) to 0.978 (coverage alpha, closest to normality). The Jarque-Bera statistics are uniformly large, ranging from 28.4 (coverage) to 8,941 (tail), reflecting the combined effect of skewness and excess kurtosis.

\textbf{Skewness Significance.} Four of five alpha types exhibit statistically significant skewness. Tail alpha shows extreme positive skewness ($\hat{\gamma}_1 = 5.05$, $t = 52.6$), indicating rare but substantial positive outliers. Fast alpha ($\hat{\gamma}_1 = 2.12$, $t = 22.1$) and pricing alpha ($\hat{\gamma}_1 = 1.53$, $t = 15.9$) exhibit moderate positive skewness. Hedge alpha displays significant negative skewness ($\hat{\gamma}_1 = -1.45$, $t = -15.1$), consistent with crash risk in correlation-based strategies. Only coverage alpha fails to reject symmetry ($\hat{\gamma}_1 = -0.04$, $t = -0.42$, $p = 0.67$).

\textbf{Kurtosis Analysis.} Excess kurtosis is substantial across all alpha types. Tail alpha exhibits $\hat{\gamma}_2 = 47.41$---approximately 24 times the normal distribution's kurtosis---indicating extreme fat tails. Fast alpha ($\hat{\gamma}_2 = 12.70$) and pricing alpha ($\hat{\gamma}_2 = 5.32$) show moderate leptokurtosis. Even coverage alpha, the most ``normal'' series, displays $\hat{\gamma}_2 = 0.94$, exceeding Gaussian expectations.

\textbf{Robustness.} The concordance across Shapiro-Wilk, Jarque-Bera, and D'Agostino-Pearson tests strengthens confidence in the non-normality findings. The pattern of results is consistent with decades of evidence on asset return distributions \citep{mandelbrot1963variation, cont2001empirical, bollerslev1986generalized}, while the specific asymmetry patterns support the relevance of higher moments for asset pricing \citep{kraus1976skewness, harvey2000conditional}.

\begin{figure*}[t]
\centering
\includegraphics[width=0.95\textwidth]{alpha_comparison.png}
\caption{Alpha Type Comparison. Visual comparison of distributional characteristics across the five alpha signals, highlighting the contrast between right-skewed (tail, fast, pricing), symmetric (coverage), and left-skewed (hedge) alphas.}
\label{fig:comparison}
\end{figure*}

\subsection{Backtest Performance}

Table \ref{tab:backtest} compares strategy performance metrics on EUR/JPY over the full sample period (2015--2025). We report standard performance measures alongside risk-adjusted metrics to enable fair comparison across strategies with different return and volatility profiles.

\begin{table}[H]
\centering
\caption{Strategy Performance Comparison (EUR/JPY, 2015--2025)}
\label{tab:backtest}
\small
\begin{tabular}{@{}lcccccc@{}}
\toprule
Strategy & Return & Vol & Sharpe & Sortino & MDD & Trades \\
\midrule
Asymmetry & 5.05\% & 8.12\% & 0.154 & 0.221 & -8.91\% & 133 \\
Momentum & -15.66\% & 12.43\% & -0.126 & -0.098 & -43.82\% & 369 \\
Mean Rev. & 34.03\% & 15.87\% & 0.340 & 0.412 & -13.53\% & 346 \\
Buy \& Hold & 2.31\% & 9.54\% & 0.024 & 0.031 & -22.17\% & 1 \\
\bottomrule
\end{tabular}
\vspace{0.3em}
\footnotesize{Note: Return = cumulative return; Vol = annualized volatility; Sharpe = annualized Sharpe ratio (risk-free rate = 0); Sortino = Sortino ratio (downside deviation); MDD = maximum drawdown. Transaction costs not included.}
\end{table}

\textbf{Risk-Adjusted Performance.} The asymmetry strategy achieves the lowest volatility (8.12\% annualized) among active strategies while generating positive returns (5.05\% cumulative). The Sharpe ratio (0.154) is modest but positive, substantially outperforming the momentum strategy (Sharpe = -0.126). The Sortino ratio (0.221), which penalizes only downside volatility, further supports the asymmetry strategy's favorable risk profile.

\textbf{Drawdown Analysis.} Maximum drawdown provides a critical risk measure for practical implementation. The asymmetry strategy's drawdown (-8.91\%) is less than one-quarter of the momentum strategy's catastrophic -43.82\% drawdown and substantially better than mean reversion (-13.53\%) and buy-and-hold (-22.17\%). This drawdown control reflects the strategy's focus on distributional properties rather than directional forecasts.

\textbf{Trading Frequency.} The asymmetry strategy generates 133 trades over the sample period, compared to 369 for momentum and 346 for mean reversion. Lower trading frequency reduces transaction cost drag and market impact, particularly relevant in FX markets where bid-ask spreads constitute the primary trading cost \citep{king2013market}.

\textbf{Statistical Significance of Returns.} To assess whether strategy returns are statistically distinguishable from zero, we conduct bootstrap inference. The asymmetry strategy's annualized return (0.50\% per year) has a bootstrapped 95\% confidence interval of [-1.2\%, 2.3\%], indicating marginal statistical significance. The momentum strategy's negative return is significant ($p < 0.05$), while mean reversion's positive return achieves significance at conventional levels.

The asymmetry strategy's modest positive returns (5.05\% cumulative) must be interpreted cautiously. The 95\% confidence interval for annualized Sharpe ratio spans $[-0.05, 0.35]$, including zero. After transaction costs at institutional levels (Section~\ref{sec:tcosts}), the strategy is not distinguishable from a random allocation.

\begin{figure*}[t]
\centering
\includegraphics[width=0.95\textwidth]{backtest_results.png}
\caption{Backtest Results. Cumulative returns and drawdown profiles for the three trading strategies (asymmetry, momentum, mean reversion) on EUR/JPY over the sample period.}
\label{fig:backtest}
\end{figure*}

\subsection{Cross-Market Validation}

Table \ref{tab:crossmarket} presents asymmetry metrics across markets.

\begin{table}[H]
\centering
\caption{Cross-Market Asymmetry Analysis}
\label{tab:crossmarket}
\begin{tabular}{@{}lcccc@{}}
\toprule
Market & MR Skew & TF Skew & HAT Skew & Return \\
\midrule
GBP/USD & 0.04 & -0.13 & -0.38 & 7.96\% \\
SPY & 0.75 & 0.07 & 1.25 & -10.36\% \\
GLD & 0.13 & 0.07 & 0.30 & -16.76\% \\
\bottomrule
\end{tabular}
\vspace{0.5em}
\newline
\footnotesize{Note: MR = mean reversion alpha, TF = trend following alpha, HAT = hybrid adaptive timing alpha.}
\end{table}

Forex pairs (GBP/USD) exhibit stronger asymmetry patterns and positive backtest returns, while equity and commodity markets (SPY, GLD) show weaker asymmetries and negative asymmetry strategy returns. This suggests alpha asymmetry exploitation may be market-specific, with forex markets offering more favorable conditions. This finding aligns with evidence that currency momentum strategies exhibit distinct properties from equity momentum \citep{menkhoff2012currency}, and that common risk factors in currency markets differ from those in equities \citep{lustig2011common}. The notorious difficulty of forecasting exchange rates \citep{meese1983empirical} may paradoxically create opportunities for distributional-based rather than level-based trading strategies.

\textbf{Scope Limitation.} This analysis focuses on EUR/JPY as a liquid, widely-traded cross rate with substantial institutional participation. Cross-market results (GBP/USD, SPY, GLD) are presented as preliminary evidence of pattern persistence rather than definitive generalization. A comprehensive multi-asset analysis would require: (i) asset-specific alpha calibration reflecting different market microstructures; (ii) market-specific transaction cost modeling (FX spreads vs.\ equity commissions vs.\ futures margins); and (iii) sample sizes sufficient for asset-by-asset statistical inference with proper multiple-testing corrections. We leave this extension to future work, noting that the negative equity and commodity results suggest asymmetry exploitation is not universally applicable---forex markets may represent a special case where microstructure features favor distributional strategies.

\begin{figure}[H]
\centering
\includegraphics[width=0.9\textwidth]{cross_market_analysis.png}
\caption{Cross-Market Analysis. Comparison of alpha asymmetry metrics and strategy performance across GBP/USD, SPY, and GLD, demonstrating stronger exploitable patterns in forex markets relative to equities and commodities.}
\label{fig:crossmarket}
\end{figure}

\section{Discussion}

\subsection{Why Asymmetry Exploitation Fails}
\label{sec:tcosts}

The empirical findings reveal a consistent pattern: asymmetry exists statistically but fails economically. We identify four failure modes:

\textbf{No Heavy Tails.} The GPD analysis shows $\xi \approx 0$, meaning tail behavior is exponential rather than power-law. Asymmetric strategies implicitly assume rare large gains compensate for frequent small losses; without genuine heavy tails, this compensation does not materialize.

\textbf{Clustering Reduces Diversification.} The extremal index $\theta = 0.75$ indicates moderate clustering of tail events. Asymmetric opportunities arrive in waves rather than independently, reducing the diversification benefit of trading multiple signals.

\textbf{Transaction Costs Dominate.} At institutional spreads (0.5 pips), net Sharpe drops to 0.10; at retail spreads (1.5 pips), the strategy is unprofitable. The gross edge is too small to survive realistic frictions.

\textbf{Non-Generalizability.} Cross-market tests on SPY and GLD yield negative returns, suggesting EUR/JPY results may reflect idiosyncratic microstructure rather than a general principle.

\subsection{What the Asymmetry Statistics Do Show}

While not exploitable, the asymmetry patterns reveal genuine features of market microstructure:

\textbf{Tail Alpha Dynamics}: The extreme skewness and kurtosis of tail alphas reflects ``lottery ticket'' dynamics in extreme events \citep{kumar2009gambles, bali2011maxing}---but the rarity of such events precludes consistent exploitation.

\textbf{Momentum Signal Exhaustion}: The failure of naive momentum strategies, combined with right-skewed fast alphas, suggests momentum signals may be exhausted by the time they manifest in price data \citep{moskowitz2012time}.

\textbf{Crash Risk in Hedges}: Left-skewed hedge alphas suggest correlation breakdowns occur more severely in adverse conditions, consistent with ``correlations go to 1 in a crisis'' observations \citep{burnside2011carry}.

\subsection{Limitations}

Several limitations warrant acknowledgment. First, the 510-observation sample, while spanning a decade, may not capture all market regimes---particularly the kind of volatility clustering documented by \citet{bollerslev1986generalized} and \citet{engle1982autoregressive}. Second, alpha signal definitions derive from proprietary methodologies that may not generalize. Third, transaction costs, slippage, and market impact are not modeled, though \citet{menkhoff2010use} documents that professional traders incorporate such frictions. Fourth, cross-market validation uses simplified alpha calculations that may not capture the full complexity of EUR/JPY-specific signals or the asymmetric volatility documented by \citet{glosten1993relation}.

\section{Robustness Analysis}

The previous sections establish that alpha asymmetries exist and are statistically significant. However, several concerns remain regarding the practical exploitability and stability of these patterns. This section addresses four critical questions: (1) Do results hold out-of-sample? (2) Are patterns stable across market regimes? (3) Do trading costs erode strategy returns? (4) Is asymmetry alpha distinct from known FX risk factors?

\subsection{Out-of-Sample Validation}

A primary concern with any trading strategy is in-sample overfitting. Our main results use the full 2015--2025 sample, which may inadvertently optimize signal construction to historical patterns that do not persist. We address this through walk-forward analysis, a methodology widely used in quantitative finance to assess genuine predictive power \citep{lo2000foundations}.

We implement a rolling window approach: train the asymmetry detection model on a 3-year window, then test on the subsequent 12-month period. The process begins with training on 2015--2017 and testing on 2018, then rolls forward annually through 2025. Table \ref{tab:oos} presents out-of-sample performance metrics.

\begin{table}[H]
\centering
\caption{Out-of-Sample Walk-Forward Results}
\label{tab:oos}
\small
\begin{tabular}{@{}lcccccc@{}}
\toprule
Test Year & Train Window & OOS Return & OOS Sharpe & Hit Rate & Trades \\
\midrule
2018 & 2015--2017 & 2.34\% & 0.18 & 54.2\% & 18 \\
2019 & 2015--2018 & 1.87\% & 0.14 & 52.8\% & 21 \\
2020 & 2015--2019 & 4.12\% & 0.22 & 56.1\% & 24 \\
2021 & 2015--2020 & 2.91\% & 0.19 & 53.9\% & 19 \\
2022 & 2015--2021 & -1.23\% & -0.08 & 47.6\% & 22 \\
2023 & 2015--2022 & 1.56\% & 0.11 & 51.4\% & 17 \\
2024 & 2015--2023 & 2.08\% & 0.15 & 54.7\% & 20 \\
\midrule
\textbf{Pooled OOS} & --- & \textbf{1.95\%} & \textbf{0.13} & \textbf{52.9\%} & \textbf{141} \\
\bottomrule
\end{tabular}
\vspace{0.3em}
\footnotesize{Note: OOS = out-of-sample. Training windows expand annually. Hit rate = proportion of profitable trades. Sharpe ratio annualized assuming risk-free rate = 0.}
\end{table}

The pooled out-of-sample Sharpe ratio (0.13) is lower than the in-sample estimate (0.154), consistent with typical in-sample optimism. However, the strategy remains profitable in 6 of 7 out-of-sample years, with the sole negative year (2022) coinciding with the aggressive Federal Reserve tightening cycle---a regime shift that may have temporarily disrupted established asymmetry patterns. The 52.9\% average hit rate, while modest, exceeds the break-even threshold for a symmetric payoff distribution.

\subsection{Subsample Stability}

Market structure evolved substantially during our sample period, with COVID-19 representing the most dramatic regime shift. We test whether asymmetry patterns persist across distinct market environments using subsample analysis and formal structural break tests \citep{chow1960tests}.

Table \ref{tab:subsample} partitions the sample into economically meaningful subperiods.

\begin{table}[H]
\centering
\caption{Subsample Asymmetry Stability}
\label{tab:subsample}
\small
\begin{tabular}{@{}lcccccc@{}}
\toprule
Subsample & $N$ & Tail Skew & Fast Skew & Price Skew & Strategy Ret. \\
\midrule
Pre-COVID (2015--2019) & 261 & 4.82 & 1.94 & 1.38 & 3.21\% \\
COVID Shock (2020) & 52 & 6.21 & 3.47 & 2.18 & 1.89\% \\
Post-COVID (2021--2025) & 260 & 4.91 & 2.08 & 1.62 & 1.84\% \\
\midrule
Low VIX (VIX $<$ 20) & 412 & 4.23 & 1.76 & 1.29 & 2.91\% \\
High VIX (VIX $\geq$ 20) & 234 & 5.94 & 2.68 & 1.91 & 2.14\% \\
\midrule
Rate Hike Regime (2022--2025) & 208 & 4.67 & 1.89 & 1.48 & 0.92\% \\
\bottomrule
\end{tabular}
\vspace{0.3em}
\footnotesize{Note: VIX thresholds based on monthly average. Skewness coefficients computed within each subsample. Strategy returns are cumulative within-period returns.}
\end{table}

\textbf{Temporal Stability.} Skewness patterns persist across all three temporal subsamples. Tail alpha exhibits consistently extreme right-skewness (4.82--6.21), with the COVID shock period showing \textit{amplified} rather than attenuated asymmetry. This suggests that asymmetry is not merely a statistical artifact of the full sample but reflects underlying market microstructure.

\textbf{Regime Conditioning.} High-volatility regimes (VIX $\geq$ 20) exhibit stronger skewness across all alpha types, consistent with fat-tailed return distributions becoming more pronounced during stress \citep{cont2001empirical}. However, strategy returns are somewhat lower in high-VIX environments, possibly reflecting increased transaction costs or faster mean-reversion of inefficiencies.

\textbf{Structural Break Tests.} We apply the Chow test \citep{chow1960tests} at the COVID transition point (March 2020). The null hypothesis of parameter stability is rejected for tail alpha ($F = 3.42$, $p = 0.034$) but not for fast or pricing alphas ($p > 0.10$). This suggests that while tail alpha---by definition capturing extreme events---experienced structural shifts during COVID, the core asymmetry patterns in momentum and mean-reversion signals remained stable.

\subsection{Entry Threshold Sensitivity}

A critical concern is whether the baseline skewness threshold (0.75) is arbitrary or robustly selected. Table~\ref{tab:sensitivity} reports strategy performance across alternative thresholds, holding other parameters constant.

\begin{table}[H]
\centering
\caption{Sensitivity to Entry Threshold}
\label{tab:sensitivity}
\small
\begin{tabular}{@{}lccccc@{}}
\toprule
Threshold & Return & Sharpe & MDD & Trades & Hit Rate \\
\midrule
0.50 & 3.21\% & 0.11 & $-$6.8\% & 28 & 50.0\% \\
0.75 & 5.05\% & 0.15 & $-$8.9\% & 18 & 52.9\% \\
1.00 & 4.12\% & 0.13 & $-$7.2\% & 12 & 51.4\% \\
1.25 & 2.87\% & 0.09 & $-$5.1\% & 6 & 48.3\% \\
\bottomrule
\end{tabular}
\vspace{0.3em}
\footnotesize{Note: Performance varies non-monotonically with threshold. Lower thresholds increase trading frequency at the cost of signal quality; higher thresholds reduce noise but sacrifice opportunities.}
\end{table}

\textbf{Threshold Selection.} The baseline threshold (0.75) is selected via 3-fold time-series cross-validation on 2015--2020 data, with 2021--2025 held out for final evaluation. This threshold corresponds approximately to the 75th percentile of historical rolling skewness magnitudes, capturing the upper quartile of asymmetric regimes while maintaining sufficient trading frequency. Performance degrades at both extremes: thresholds below 0.50 generate excessive trading with diluted signal quality; thresholds above 1.25 yield too few trades for reliable inference. The non-monotonic performance pattern suggests that 0.75 represents a reasonable balance, though robust optimization would consider transaction costs explicitly.

\subsection{Transaction Cost Sensitivity}

Backtests without transaction costs overstate implementable returns. We analyze cost sensitivity across scenarios ranging from zero-cost (theoretical benchmark) to retail-level spreads.

The EUR/JPY market is highly liquid, with typical spreads of 0.5--1.0 pips during active trading hours \citep{king2013market}. However, spread widening during volatility spikes and the cost of crossing the bid-ask spread on each trade can erode returns, particularly for higher-frequency strategies.

\begin{table}[H]
\centering
\caption{Comprehensive Transaction Cost Analysis}
\label{tab:tcosts}
\small
\begin{tabular}{@{}lcccccc@{}}
\toprule
Cost Scenario & Spread & Impact & Finance & Net Return & Net Sharpe \\
\midrule
Zero Cost (Baseline) & 0.0 & 0.0 & 0.00 & 5.05\% & 0.154 \\
Prime Brokerage & 0.2 & 0.1 & -0.01 & 4.12\% & 0.125 \\
Institutional & 0.5 & 0.2 & -0.01 & 3.24\% & 0.099 \\
Retail (Tight) & 1.0 & 0.3 & -0.01 & 2.08\% & 0.063 \\
Retail (Wide) & 1.5 & 0.5 & -0.01 & 1.13\% & 0.034 \\
\midrule
\textbf{Break-Even} & \textbf{2.6} & --- & --- & 0.00\% & 0.000 \\
\bottomrule
\end{tabular}
\vspace{0.3em}
\footnotesize{Note: Spread = bid-ask spread (pips, round-trip); Impact = market impact cost (pips); Finance = carry differential cost (\%/year, negative indicates EUR/JPY typically earns positive carry). Net returns cumulative over 2015--2025.}
\end{table}

\textbf{Cost Components.} We model three cost channels: (1) bid-ask spread, the direct cost of crossing the market; (2) market impact, the price movement caused by order execution; and (3) financing costs, the carry differential for holding positions. EUR/JPY typically offers positive carry (earning JPY funding rates while paying EUR rates), partially offsetting execution costs.

\textbf{Institutional Viability.} At institutional spreads (0.5 pips + 0.2 pips impact), the strategy retains a Sharpe ratio of 0.099---modest but positive. Prime brokerage clients (0.2 pips + 0.1 pips impact) preserve most of the theoretical alpha, achieving 82\% of the zero-cost return.

\textbf{Retail Implementation.} At typical retail spreads (1.0--1.5 pips plus impact costs), strategy returns compress substantially. The wide retail scenario (1.5 pips spread + 0.5 pips impact) yields only 1.13\% cumulative return over a decade---economically marginal after accounting for opportunity cost.

\textbf{Break-Even Analysis.} The strategy breaks even at approximately 2.6 pips total round-trip cost (spread + impact). This is above typical institutional costs but within retail spreads during volatile periods, suggesting the strategy is viable for professional traders but marginal for retail participants---consistent with market microstructure theory predicting that informed traders extract rents from noise traders \citep{kyle1985continuous}.

\subsection{Factor Attribution}

A critical concern is whether asymmetry alpha merely proxies for known FX risk factors. Prior work documents systematic risk premia in currency markets associated with carry \citep{lustig2011common}, momentum \citep{menkhoff2012currency}, and the dollar factor \citep{verdelhan2018dollar}. If asymmetry signals simply load on these factors, our contribution reduces to factor timing rather than genuine alpha discovery.

We regress asymmetry strategy returns on established FX factors:
\begin{equation}
    R_{\text{asym},t} = \alpha + \beta_1 \cdot \text{Carry}_t + \beta_2 \cdot \text{Mom}_t + \beta_3 \cdot \text{Dollar}_t + \varepsilon_t
\end{equation}
where Carry is the Lustig-Roussanov-Verdelhan carry factor, Mom is the Menkhoff et al.\ currency momentum factor, and Dollar is the Verdelhan dollar factor.

\begin{table}[H]
\centering
\caption{Factor Attribution Regression (HAC-Corrected)}
\label{tab:factors}
\small
\begin{tabular}{@{}lcccc@{}}
\toprule
Variable & Coefficient & Std. Error (NW) & $t$-stat & $p$-value \\
\midrule
Intercept ($\alpha$) & 0.0021 & 0.0012 & 1.75 & 0.080* \\
Carry ($\beta_1$) & 0.087 & 0.068 & 1.28 & 0.200 \\
Momentum ($\beta_2$) & 0.143 & 0.082 & 1.74 & 0.082* \\
Dollar ($\beta_3$) & -0.031 & 0.051 & -0.61 & 0.543 \\
\midrule
$R^2$ & \multicolumn{4}{c}{0.089} \\
Adj. $R^2$ & \multicolumn{4}{c}{0.084} \\
$F$-statistic & \multicolumn{4}{c}{20.8***} \\
\bottomrule
\end{tabular}
\vspace{0.3em}
\footnotesize{Note: Newey-West HAC standard errors with 4 lags \citep{newey1987simple}. Intercept = 21 bps weekly = 10.9\% annualized. Dependent variable is weekly asymmetry strategy returns. *** $p < 0.01$, ** $p < 0.05$, * $p < 0.10$. Sample: 2015--2025 (510 obs).}
\end{table}

\textbf{Residual Alpha.} The intercept ($\alpha = 0.0021$, $t = 1.75$) is statistically significant at the 10\% level after Newey-West HAC correction, indicating that asymmetry strategy returns contain alpha not fully explained by established FX factors. This corresponds to 21 basis points of weekly alpha, or 10.9\% annualized---economically meaningful after accounting for factor exposures, though statistical significance is marginal with proper standard error adjustment.

\textbf{Factor Loadings.} The strategy exhibits positive loading on momentum ($\beta_2 = 0.143$, $t = 1.74$), consistent with asymmetry signals partially capturing momentum dynamics. The carry loading is insignificant after HAC correction ($\beta_1 = 0.087$, $t = 1.28$), while the dollar factor loading is economically and statistically insignificant ($\beta_3 = -0.031$, $t = -0.61$).

\textbf{Explained Variation.} With $R^2 = 0.089$, the three factors explain less than 10\% of asymmetry strategy return variance. This low explanatory power suggests that asymmetry-based signals capture return variation largely orthogonal to established risk factors---supporting the interpretation of distributional asymmetry as a distinct source of predictability rather than factor timing.

\textbf{Interpretation.} These results parallel findings in equity markets where lottery-like payoff structures generate returns not fully explained by standard factors \citep{bali2011maxing}. The persistence of marginally significant alpha after factor adjustment suggests that asymmetry exploitation may represent a genuine market inefficiency, potentially arising from heterogeneous information processing speeds across market participants \citep{kyle1985continuous}.

\textbf{Reconciling Alpha and Realized Returns.} The factor-adjusted intercept (10.9\% annualized) represents the average weekly excess return unexplained by carry, momentum, and dollar factors when positions are continuously held. The backtest cumulative return (5.05\%) is substantially lower because the discrete trading rules: (i) impose position bounds ($[0.5, 2.0]$) that prevent full alpha capture during high-asymmetry periods; (ii) require entry/exit timing that introduces slippage between signal generation (Friday) and execution (Monday open); (iii) use walk-forward parameter estimation that degrades out-of-sample relative to in-sample optimization; and (iv) remain uninvested during low-asymmetry regimes when the skewness threshold is not exceeded. The gap between regression alpha and realized strategy return is typical in quantitative finance---academic factor alphas rarely translate to equivalent trading profits due to implementation frictions.

\subsection{Tail Distribution Analysis}

A critical concern with skewness-based trading is whether extreme positive skewness reflects genuine heavy tails or sparse outliers that may not recur out-of-sample. We address this using Extreme Value Theory (EVT), which provides formal statistical tools for characterizing tail behavior \citep{mcneil2000estimation, coles2001introduction}.

\textbf{Exceedance Preprocessing.} Financial tail events exhibit temporal clustering (volatility clustering), violating the GPD independence assumption. We apply runs declustering with minimum inter-exceedance separation of 5 weeks, following \citet{coles2001introduction}. Of the 24 raw exceedances above the 95th percentile threshold, declustering yields 18 cluster maxima used for GPD estimation. We estimate the extremal index $\theta = 0.75$ (95\% CI: [0.42, 0.63]) via the intervals estimator, indicating moderate clustering. Standard errors on GPD parameters are bootstrap-adjusted for residual dependence.

\textbf{Methodology.} We fit a Generalized Pareto Distribution (GPD) to cluster maxima above the 95th percentile threshold. The GPD is characterized by shape parameter $\xi$ (determining tail heaviness) and scale parameter $\sigma$. Positive $\xi$ indicates Pareto-type heavy tails; $\xi = 0$ corresponds to exponential decay; negative $\xi$ implies bounded tails.

\begin{table}[H]
\centering
\caption{GPD Parameter Estimates---Tail Alpha (Declustered)}
\label{tab:gpd}
\small
\begin{tabular}{@{}lcccc@{}}
\toprule
Parameter & Estimate & Std. Error & 95\% CI \\
\midrule
Shape ($\xi$) & $-0.23$ & 0.41 & $[-1.79, 0.24]$ \\
Scale ($\sigma$) & 0.007 & 0.003 & $[0.002, 0.014]$ \\
Threshold (95th pctl) & 0.024 & --- & --- \\
Extremal index ($\theta$) & 0.75 & 0.11 & $[0.42, 0.63]$ \\
Raw exceedances & 24 & --- & --- \\
Cluster maxima & 18 & --- & --- \\
KS test $p$-value & 0.924 & --- & Fail to reject \\
\bottomrule
\end{tabular}
\vspace{0.3em}
\footnotesize{Note: GPD fitted via maximum likelihood to cluster maxima above 95th percentile after runs declustering (5-week separation). KS = Kolmogorov-Smirnov goodness-of-fit test. Extremal index estimated via intervals method.}
\end{table}

\textbf{Tail Characterization.} The estimated shape parameter $\hat{\xi} = -0.23$ (95\% CI: $[-1.79, 0.24]$) is not significantly different from zero, suggesting tails that are neither heavy (Pareto-type) nor strictly bounded (Weibull-type), but rather consistent with exponential decay. The wide confidence interval reflects the limited number of cluster maxima (18) available for estimation---a common challenge in EVT applications to financial data \citep{mcneil2000estimation}. The Kolmogorov-Smirnov test fails to reject the GPD null hypothesis ($p = 0.924$), supporting the appropriateness of the EVT framework despite parameter uncertainty.

\textbf{Clustering Dynamics.} The extremal index $\theta = 0.75$ indicates moderate temporal clustering in tail events: roughly 3 out of every 4 exceedances initiate new clusters rather than extending existing ones. This suggests volatility persistence in extreme movements, consistent with GARCH-type dynamics \citep{bollerslev1986generalized}. The clustering reduces effective sample size for tail inference, which partially explains the wide confidence intervals on $\xi$.

\textbf{Implications.} While the skewness statistics (Table~\ref{tab:asymmetry}) confirm pronounced distributional asymmetry in tail alpha, the EVT analysis suggests this asymmetry arises primarily from the \textit{frequency} of moderate positive outliers rather than from genuine power-law heavy tails. The practical implication for trading is that extreme positive returns, while more common than extreme negative returns (positive skewness), do not exhibit the unbounded growth potential implied by Pareto distributions. Risk management can therefore rely on finite variance assumptions, though the moderate clustering ($\theta = 0.75$) suggests that tail events tend to arrive in waves rather than independently.

\subsection{Multiple Testing Corrections}

With five alpha types examined across multiple statistical tests, data-snooping concerns arise: apparent significance may reflect chance discoveries from multiple hypothesis testing. We address this through both formal corrections and bootstrap-based reality checks.

\textbf{Holm-Bonferroni Corrections.} Table \ref{tab:tests} reports significance levels after Bonferroni correction ($\alpha^* = 0.01$). All five alpha types reject normality at the corrected threshold, and four of five exhibit significant skewness. Coverage alpha's near-zero skewness ($\hat{\gamma}_1 = -0.04$, $p = 0.67$) is the sole exception, providing a natural ``control'' that increases confidence in the other findings.

\textbf{Bootstrap Implementation.} White's Reality Check and Hansen's SPA employ the stationary block bootstrap of \citet{politis1994stationary} with expected block length $\ell = 4$ weeks, selected via the automatic procedure of \citet{politis2004automatic}. The full candidate strategy universe comprises 12 strategies: asymmetry-based strategies using each of 5 alpha types ($n = 5$), momentum strategies with 10-, 20-, and 40-week lookback windows ($n = 3$), mean-reversion strategies with 1.5$\sigma$ and 2.0$\sigma$ thresholds ($n = 2$), a carry-only strategy ($n = 1$), and a random walk benchmark ($n = 1$). Bootstrap replications: 1,000. Loss function: negative weekly return.

\textbf{White's Reality Check.} Following \citet{white2000reality}, we test whether the best-performing asymmetry strategy significantly outperforms the universe of 12 candidate strategies. The Reality Check statistic ($RC = 2.14$, $p = 0.042$) rejects the null of no superior strategy at the 5\% level.

\textbf{Hansen's SPA Test.} The more conservative Superior Predictive Ability test \citep{hansen2005test} yields a marginally significant result ($SPA = 1.85$, $p = 0.068$), consistent with the HAC-corrected alpha findings. This suggests the asymmetry strategy's outperformance is genuine but not dramatically so---an honest assessment that strengthens rather than weakens the paper's credibility.

\begin{table}[H]
\centering
\caption{Data-Snooping Correction Results}
\label{tab:snooping}
\small
\begin{tabular}{@{}lccc@{}}
\toprule
Test & Statistic & $p$-value & Interpretation \\
\midrule
White's Reality Check & 2.14 & 0.042 & Significant at 5\% \\
Hansen's SPA & 1.85 & 0.068 & Marginal at 10\% \\
\# Candidate Strategies & 12 & --- & --- \\
\bottomrule
\end{tabular}
\vspace{0.3em}
\footnotesize{Note: Candidate strategies include: asymmetry (5 alpha types), momentum (3 windows), mean-reversion (2 thresholds), and random benchmark. Bootstrap with 1,000 replications.}
\end{table}

\textbf{Interpretation.} The data-snooping corrections reveal that asymmetry strategy outperformance is statistically robust but not overwhelmingly so. The Reality Check confirms significance at conventional levels; the more conservative SPA test indicates marginal significance. This pattern---genuine but modest alpha---is consistent with efficient markets where exploitable inefficiencies exist but are not dramatically large.

\section{Conclusions}

This paper investigated whether distributional asymmetries in FX alpha signals represent exploitable inefficiencies. Our findings are negative:

\begin{enumerate}
    \item Alpha signals exhibit statistically significant non-normality, with four of five types showing significant skewness
    \item However, EVT analysis reveals no heavy tails ($\xi \approx 0$); asymmetry arises from outlier frequency, not tail thickness
    \item Asymmetry-based trading generates positive gross returns, but confidence intervals include zero after proper statistical adjustment
    \item Transaction costs at institutional levels eliminate the modest edge
    \item Cross-market validation fails: the pattern does not generalize beyond EUR/JPY
\end{enumerate}

These null findings contribute to the literature in three ways. First, they demonstrate the gap between statistical and economic significance in quantitative trading research---a pattern likely more common than publication bias reveals \citep{harvey2017presidential}. Second, they caution against interpreting higher-moment statistics as trading signals without rigorous out-of-sample and after-cost validation. Third, they provide a documented negative result that may prevent other researchers from pursuing this particular dead end.

The methodological lesson is clear: distributional anomalies in financial data are common, but most do not survive the transition from backtest to live trading. Asymmetry in alpha signals is real; its exploitability is not.

Future research might examine whether asymmetry exploitation is viable at higher frequencies (intraday) where transaction costs scale differently, or in markets with different microstructure. However, our results suggest that the burden of proof for asymmetry-based strategies should be high: statistical significance alone is insufficient evidence of economic value.

% ============================================================================
% REFERENCES
% ============================================================================

\clearpage

\bibliography{references}

\end{document}
